\section{Representation of edge laws}
\label{sec:a-laws}

The structural reduction of \cref{sec:a-block-rank} uses monotonicity,
whereas its exact algorithm uses a polynomial representation. We first
separate these roles for continuous piecewise-polynomial laws with
independent uncertain coefficients. We then show why a fixed fractional
power changes exact arc comparison, even on one cycle. Finally, positive
rational approximations recover an algorithm polynomial in accuracy bits
for a fixed finite family of fractional powers. All optimization domains
in this section are unfiltered product domains: operating capacities are
tested after optimization, not imposed as scenario constraints.

\subsection{Dense polynomial laws and their admissibility}
\label{subsec:a-law-model}
An edge has the law
\begin{equation}
 g_e(x,\theta_e)=f_{e0}(x)+\sum_{j=1}^{h_e}\theta_{ej}f_{ej}(x),
 \qquad \theta_{ej}\in[\underline\theta_{ej},\overline\theta_{ej}].
 \label{eq:a-law-affine}
\end{equation}
The intervals are finite and rational, and all their coordinates are
independent, including across edges. Every basis function is continuous
on $\R$, vanishes at zero, and is specified by a finite list of rational
breakpoints and rational polynomial pieces, including its two unbounded
outer intervals. Polynomials are densely encoded: coefficients through
the stated degree are listed, including zeros. Degrees, piece counts,
and the $h_e$ may grow with input size. The complete law must be strictly
increasing for every allowed coefficient vector. Individual bases need
not be increasing, and coefficient intervals need not be positive.
There is no derivative-matching condition at breakpoints.

This affine notation will also be useful for explicitly stated correlated
models; the results proved here use independent boxes. In particular, a
parameter shared by different edges is not a coordinate $\theta_{ej}$ in
this model. Fixed laws correspond to $h_e=0$. Positive uncertain mixtures
of $x$, $x|x|$, $x^3$, and $\max\{x-a,0\}$ with rational $a\ge0$ give
examples whenever the complete law is strictly increasing. Signed even
integer powers use two polynomial pieces.

\begin{proposition}[Validation of the law family]
\label{prop:a-law-validation}
All the representation and admissibility conditions above can be checked
in polynomial bit time in the dense rational input size.
\end{proposition}
\begin{proof}
Check breakpoint order and discard empty pieces. Rational evaluation
checks continuity at every breakpoint and the value at zero of every
basis. Refine each edge's pieces by the union of its basis breakpoints.
The total number of refined descriptions is polynomial in the explicit
input size; this operation does not enumerate combinations of pieces.
On an open refined interval $J$, the minimum derivative over the box is
\begin{equation}
 f'_{e0}(x)+\sum_j\min\{
 \underline\theta_{ej}f'_{ej}(x),
 \overline\theta_{ej}f'_{ej}(x)\}.
 \label{eq:a-law-min-derivative}
\end{equation}
Discard identically zero $f'_{ej}$ when forming a root partition. Isolate
and order the real roots in $J$ of the product of all remaining $f'_{ej}$.
Its degree and coefficient bits are polynomial in dense input length.
Between consecutive roots each endpoint choice in
\cref{eq:a-law-min-derivative} is fixed. Univariate polynomial sign tests
check nonnegativity of the resulting rational polynomial on each such
interval, including unbounded intervals. At partition roots use continuity
of this minimum of finitely many continuous polynomials. One-sided limits
suffice at original breakpoints. These tests are polynomial-time instances
of the real-algebraic primitives in \cref{subsec:a-pre-computation}.
They prove, or refute, that every complete law is nondecreasing.

For each original refined interval $J$, also test the rational linear
system that sets every coefficient of
$f'_{e0}+\sum_j\theta_{ej}f'_{ej}$ to zero, together with the coefficient
box. Rational LP decides its feasibility in polynomial time. A feasible
system supplies an allowed law constant on $J$, which fails strict
increase. Conversely, a continuous nondecreasing function failing strict
increase is constant on a nontrivial interval. Some open subinterval
lies in one refined piece, and a polynomial derivative vanishing there
vanishes identically on that piece. Thus one of these LPs is feasible.
Reject exactly in this case. There are polynomially many LPs and sign
tests. No field containing all isolated roots simultaneously is needed.
\end{proof}

The strictness test is essential. The family $x+\theta|x|$ with
$\theta\in[-1,1]$ is nondecreasing but has a flat half-line at either
endpoint; the smaller box $[-1/2,1/2]$ is admissible. In contrast,
$(1-\theta)x^3+\theta x$, $\theta\in[0,1]$, is admissible despite a zero
derivative at $x=0,\theta=0$. Isolated derivative zeros are permitted.

\subsection{Structural transfer, dense degrees, and coefficient recovery}
\begin{lemma}[Law-independent nomination faces]
\label{lem:a-law-faces}
For any fixed continuous strictly increasing laws vanishing at zero,
the graph-defined face family of \cref{lem:a-blk-faces} contains a
maximizer of a prescribed terminal drop. Its size and free-coordinate
bounds are unchanged. The same family contains a joint maximizer when
the laws are jointly continuous in flow and parameters in a compact parameter domain and
remain admissible throughout that domain.
\end{lemma}
\begin{proof}
Existence follows from \cref{thm:a-pre-existence}. Acyclic physical flows
give $|x_e|\le B:=\sum_v\max\{|\ell_v|,|u_v|\}$, as in the proof of
\cref{lem:a-blk-faces}. On this flow interval, compactness and continuity
of the laws bound every normalized potential by a sum of edge-drop
bounds along a path. If inputs converge, every subsequential state limit
satisfies the limiting equations; uniqueness identifies it. This proves
continuity of the state and attainment on the compact input domain.

To handle corners, fix a nonnegative smooth kernel $\kappa$ supported
on $[-1,1]$ with integral one, and put
\begin{align*}
 S_\rho f(x)&=\int\kappa(t)f(x-\rho t)\,dt
                  -\int\kappa(t)f(-\rho t)\,dt,\\
 g_{e,\rho}(x)&=S_\rho g_e(x)+\rho x,\qquad 0<\rho\le1.
\end{align*}
Convolution preserves increase, centering enforces value zero at zero,
and $g'_{e,\rho}\ge\rho$. The smoothed laws converge uniformly on the
flow interval to the original laws; for a compact continuous parameter
family this convergence is uniform in the parameters too. The preceding
bounded-state subsequence argument gives uniform convergence of physical
states. Centering is what preserves the sign of flow and its acyclic
bound in this argument.

For these smooth laws the differentiated network has positive diagonal
resistances $g'_{e,\rho}(x_e)$. The proofs of
\cref{lem:a-blk-order,lem:a-blk-unimodal} therefore apply: the unit
adjoint has ordered block ranges, and positive forcing $\delta$ at the
internal vertices of a suppressed path gives
$\jmath_i-\jmath_{i-1}=\delta$. The same threshold patterns leave at
most $8r_H-2$ free nominations in the selected nonbridge block. The
perturbing sum of normalized potentials is uniformly bounded, so its
contribution tends uniformly to zero with $\delta$. Let $\rho,\delta$
tend to zero and take a subsequence in one of the finite closed faces.
Its limit maximizes the original objective. This argument only proves
existence of an optimizer on a face; neither smoothing nor its rate is
part of the algorithm. Finally, at a joint maximizer fix its parameters
and replace its nomination by such a face maximizer. It remains jointly
optimal. The face family depends on the graph and nomination bounds,
not on the chosen laws.
\end{proof}

\begin{lemma}[Growing dense degree in the box lemma]
\label{lem:a-law-dense}
In \cref{lem:a-blk-boxlp}, fix only the core dimension and the number
of linking equations. Allow the degree bound to grow, with all core
polynomials densely encoded. Its exact optimization and recovery
conclusions still have polynomial bit complexity. The returned common
local algebraic field has polynomial degree and encoding length, not
necessarily constant degree.
\end{lemma}
\begin{proof}
The support scores there are $\eta^\top W_j(z)$, of degree at most
$D+1$ in a fixed number of variables. The number of realizable sign
vectors is $(h(D+1))^{O(d_0+k+1)}$. Listing only these vectors and forming
their guarded support sums therefore gives a polynomial-size formula.
Every sum has polynomial coefficient bits. The fixed-dimensional
quantifier-elimination bound \cref{eq:a-pre-qe-cost} is polynomial in
polynomial count, numerical degree, and coefficient bits; it does not
require degree to be fixed. Apply it also to the optimal-pair formula,
then use \cref{prim:a-pre-sample-points}. Degrees and encoding lengths
of the resulting algebraic sample are polynomial for fixed dimensions.

For recovery, all columns evaluated at this sample lie in its one
ordered algebraic field. The construction in \cref{lem:a-blk-boxlp}
enumerates polynomially many support cells in fixed dimension, then
subsets of at most $k+2$ image vertices and fixed-size affine systems.
Exact field arithmetic and sign determination are polynomial in the
field degree and coefficient bits. Fixed-size determinants and the
final polynomially many sums keep these sizes polynomial. Thus growing
local degree does not create a compositum of independently represented
roots. Dense encoding ensures $D$ is bounded by input size; sparse
binary-encoded exponents would not give this bound.
\end{proof}

For later rational recovery define basis bounds $N_{ej},M_{ej}\in\Q_{\ge0}$
on $[-B-1,B+1]$ by taking, over all pieces $\sum_i a_ix^i$, the maxima of
\[
 \sum_i|a_i|T^i,\qquad
 \sum_{i\ge1}i|a_i|T^{i-1},\qquad T=\max\{1,B+1\},
\]
respectively. Zero bounds are allowed. Put
$U_{ej}=\max\{|\underline\theta_{ej}|,|\overline\theta_{ej}|\}$ and
$M_e=M_{e0}+\sum_jU_{ej}M_{ej}$. These rationals have polynomial bit
length even for growing dense degrees.

\begin{lemma}[Pressure sensitivity for affine law boxes]
\label{lem:a-law-sensitivity}
For a simple terminal path $P_{st}$ put $C_b=\sum_{e\in P_{st}}M_e$.
Writing $F_{s,t}(b,\theta)=\pi_s-\pi_t$, we have
\begin{equation}
 |F_{s,t}(b,\theta)-F_{s,t}(c,\eta)|
 \le C_b\|b-c\|_1+
              2\sum_{e,j}N_{ej}|\theta_{ej}-\eta_{ej}|.
 \label{eq:a-law-sensitivity}
\end{equation}
\end{lemma}
\begin{proof}
Use the centered smoothing from \cref{lem:a-law-faces}. A continuous
piecewise polynomial is locally Lipschitz; its convolution derivative
is the convolution of its almost-everywhere derivative. Hence the
complete smoothed derivative lies in $[\rho,M_e+\rho]$ on physical flows.
The unit $s$--$t$ electrical adjoint has effective resistance at most
$\sum_{e\in P_{st}}(M_e+\rho)$, by testing its minimum energy against
a unit path flow. Grounding at $t$, its potential range is between zero
and that resistance, by the maximum principle. The nomination
sensitivity $dF=h^\top db$ is thus bounded by this path constant times
$\|db\|_1$.

Linearity of centered convolution preserves the affine coefficients.
At fixed $b$, differentiating the state equations and pairing with the
unit adjoint gives
\[
 \frac{\partial F_{s,t,\rho}}{\partial\theta_{ej}}
       =\jmath_e S_\rho f_{ej}(x_e).
\]
As in \cref{lem:a-blk-resistance-lipschitz}, the electrical current is
acyclic and decomposes into paths of total weight one, so
$|\jmath_e|\le1$. Also $|S_\rho f_{ej}(x_e)|\le2N_{ej}$. Integrate first
on the balanced nomination segment and then on the coefficient-box
segment. Every intermediate law is admissible. Uniform state convergence
as $\rho\downarrow0$ proves the claimed estimate.
\end{proof}

\begin{theorem}[Independent polynomial-law uncertainty]
\label{thm:a-law-polynomial}
Fix $\mathfrak r$. On a connected graph of block rank at most
$\mathfrak r$, take the law model \cref{eq:a-law-affine} and a nonempty
balanced rational nomination box. In polynomial time in the dense
input size and requested accuracy bits one can compute a rational
interval of width at most $\eps$ containing the maximum of any prescribed
terminal drop, and rational admissible nomination and coefficient inputs
whose objective is within $\eps$ of that maximum. Maximum and minimum
signed flow on any arc are exactly computable real algebraic numbers
of polynomial encoding length. Consequently unfiltered robust validation
of rational arc capacities is polynomial-time decidable, including equality.
\end{theorem}
\begin{proof}
The case $s=t$, an edgeless graph, or $B=0$ is immediate for the relevant
zero objective; rational feasible inputs are obtained by box/balance
allocation. Otherwise use \cref{lem:a-law-faces}. For each face all but
one block have fixed aggregated nominations. Since coefficients are
independent across blocks, the maximum face value is the sum of the
independently optimized block drops.

In a block, use the affine flow coordinates $X_e(z)$ of
\cref{eq:a-blk-core}. Pull back every basis breakpoint $a$ to the
affine hyperplane $X_e(z)=a$. Their arrangement has polynomial size in
the fixed core dimension. Use its closed cells, intersected with the
compact nomination and circulation polytope. Constant affine forms
are assigned their appropriate piece; at a breakpoint either adjacent
formula agrees. Thus boundary states are neither lost nor added.
On a cell every $f_{ej}(X_e(z))$ is a rational polynomial. Substitution
of a degree-$D$ polynomial into an affine form in $d$ variables has at
most $\binom{D+d}{d}$ monomials and polynomial coefficient bits when
$d$ is fixed.

For a cycle-basis row $Z_{H,i}$ and a path vector $\omega_H$, the
linking equations and objective are
\begin{align*}
 \sum_e Z_{H,ie}\Bigl(f_{e0}(X_e)
               +\sum_j\theta_{ej}f_{ej}(X_e)\Bigr)&=0,\\
 G_H(z,\theta)&=\sum_e\omega_{H,e}\Bigl(f_{e0}(X_e)
               +\sum_j\theta_{ej}f_{ej}(X_e)\Bigr).
\end{align*}
Their equivalence to the physical equations follows from orthogonality
to the cycle space, exactly as in \cref{eq:a-blk-polynomials}.
There are $r_H$ linking equations; each scalar $\theta_{ej}$ is one
interval leaf. Move the fixed terms to the right side and apply
\cref{lem:a-law-dense}. A bridge uses its conserved affine flow and no
cycle equation, so is covered as well. Compare local cell optima exactly.

Keep algebraic data for distinct blocks separate. Approximate each
block maximum by an interval of width at most $\eps/(8n)$, sum by face,
and take the maximum lower and upper endpoints. This brackets the global
maximum with width at most $\eps/8$. Selecting the face with largest
lower endpoint loses at most $\eps/8$; store its exact local optimizers.
To recover a rational input, enclose their free nomination coordinates
in rational intervals and intersect with their rational face and exact
balance. This polytope contains the algebraic nomination, so rational
LP (or box/balance allocation) returns a feasible rational nomination.
Round each coefficient in its own isolating interval intersected with
its input interval. Choose widths so that the two terms in
\cref{eq:a-law-sensitivity} are each at most $\eps/8$, ignoring zero
sensitivity constants rather than dividing by them. Only polynomially
many precision bits are required. Greedy disaggregation and arbitrary
rational coefficients off the terminal core complete the scenario.
Its new physical state need not have the sampled circulations; the
sensitivity estimate controls its pressure directly.

For an arc $e=(u,v)$, aggregate into its one block. At fixed coefficients,
strict increase of $g_e$ makes maximizing $x_e$ equivalent to maximizing
$\pi_u-\pi_v$ over nominations. Thus a joint arc maximizer lies on
one of the same faces after its coefficients are fixed. Use the affine
objective $X_e(z)$ in the box lemma. A minimum uses $-X_e$ and reversed
ordered terminals; oddness of the law is unnecessary. Exact comparison
of the polynomially many local algebraic values gives each extremum,
and rational-capacity comparison includes equality. There is no sum of
independent block values here. This conclusion does not require, or
claim, a rational near-extremal arc scenario for arbitrary polynomial laws.
\end{proof}

The general monotone-flow model and uncertainty in physical coefficients
are established. For example, \citet[Sections 3 and 4]{amann2018-deciding-robust-feasibility-and-infeasibility}
study polynomial robust feasibility, with explicit tree and cycle
reductions. \citet[Sections 1 and 3]{thurauf2025-adjustable-robust-nonlinear-network-design}
reduce robust-feasibility checking for a fixed design to polynomially many
nonlinear worst-case problems and use them in an iterative adversarial
design algorithm. That reduction is not a polynomial-time solver for each problem.
The contribution of \cref{thm:a-law-polynomial} is the fixed-block-rank
bit bound with growing dense degrees and arbitrarily many independent
coefficients, using the face reduction and fixed-dimensional support
elimination. It does not extend by analogy to uncertain breakpoints,
sparse binary exponents, or shared coefficient constraints.

\subsection{A fractional-power exact arc obstruction}
Write $\phi_q(x)=\sgn(x)|x|^q$ for $q>0$, with $\phi_q(0)=0$.
The law is algebraic for rational $q$, but that alone does not provide
a polynomial-size dense description of a sum of many such laws in a
bounded-dimensional core.

\begin{theorem}[Fractional-power arc comparison]
\label{thm:a-law-srs}
For the fixed common law $g_e(x)=\beta_e\phi_{3/2}(x)$, exact comparison
of one signed arc flow with a rational capacity is $\SRS$-hard already
on a simple cycle with fixed integral nominations and positive integral
resistances. The same holds for every fixed reduced rational exponent
$p/q>1$ whose denominator $q$ is even. It also holds for robust validation
when singleton scenario boxes are allowed.
\end{theorem}
\begin{proof}
Take positive integers $a_1,\ldots,a_k,K$ and the question
$\sum_i\sqrt{a_i}\le K$. If necessary append unit radicands and increase
$K$ accordingly to ensure $k\ge2$. On a consistently oriented simple
cycle with $k+1$ edges prescribe a base flow and resistances by
\[
 c_i=a_i,\quad\beta_i=1/a_i\quad(1\le i\le k),\qquad
 c_{k+1}=-1,\quad\beta_{k+1}=K.
\]
Set the nomination at each vertex to its outgoing base flow minus its
incoming base flow. It is integral and balanced, and every conserved
flow is $x_i=c_i+z$ for one circulation $z$. The physical equation is
\[
 h(z):=\sum_{i=1}^{k+1}\beta_i\phi_{3/2}(c_i+z)=0.
\]
It is continuous, strictly increasing, and has opposite infinite limits,
so its root $z^*$ is unique. At zero,
$h(0)=\sum_i\sqrt{a_i}-K$. Consequently, including the equality case,
\[
 \sum_i\sqrt{a_i}\le K
 \quad\Longleftrightarrow\quad h(0)\le0
 \quad\Longleftrightarrow\quad z^*\ge0
 \quad\Longleftrightarrow\quad x_{k+1}\ge-1.
\]
This is one lower capacity. Reversing that arc gives the upper capacity
$-x_{k+1}\le1$, since the law is odd. Multiplying all resistances by
$\prod_i a_i$ makes them positive integers without changing the root.
The product and every nomination have polynomial binary encoding length.

For a fixed reduced $p/q>1$ with even $q$, coprimality makes $p$ odd.
Replace the positive-edge data by
\[
 c_i=a_i^{q/2},\qquad\beta_i=a_i^{-(p-1)/2}.
\]
Then $\beta_i\phi_{p/q}(c_i)=a_i^{1/2}$. Retain the final edge and the
same capacity. The preceding equivalences hold unchanged. All exponents
in these rational data are fixed integers, so their encoding lengths
are polynomial. A common positive integer clears the resistance
denominators. Singleton boxes give the robust statement. If a convention
requires bounds on every arc, bound all other flows by the rational
acyclic-flow bound $B$; those bounds impose no further restriction.
\end{proof}

This is a lower bound for exact comparison, not $\NP$-hardness or an
additive-approximation obstruction. No converse reduction is asserted,
and the proof says nothing about fixed exponents of odd denominator.
General homogeneous laws and their arithmetic evaluation issues already
appear in \citet[Section 2.1, May 2018 manuscript]{gro2017-algorithmic-results-for-potential-based}.
The specific contribution here is the capacity predicate on a single
cycle with fixed data. A bounded source comparison does not establish
priority for all older nonlinear-circuit or hydraulic literature.

\subsection{Positive rational approximation with an explicit encoding}
\label{subsec:a-law-rational}
Positive rational approximation of fractional powers is established.
For example, the uniform half-line estimate in
\citet[Section 3.3, equation (37), Lemma 3.4, and Remark 3.1,
arXiv version]{bonito2015-numerical-approximation-of-fractional-powers}
for $\lambda^{-\alpha}$ becomes one for $t^\alpha$ on $(0,1]$ under
$\lambda=1/t$. Its positive fractions extend continuously by zero at
$t=0$. Their Section 3.2 also uses Gaussian quadrature on dyadic panels.
The following construction supplies the rational encoding and monotonicity
properties needed by the network algorithm, with an explicit proof of
its error and bit bounds.

\begin{lemma}[Positive rational scalar approximation]
\label{lem:a-law-positive-rational}
Fix a rational $0<\alpha<1$. Given a positive rational $\tau\le1$,
one can construct positive rational $a_\nu,s_\nu$ such that
\[
 R_\alpha(t)=\sum_{\nu=1}^J a_\nu\frac{t}{t+s_\nu},\qquad
 R_\alpha(1)=1,\qquad
 \sup_{0\le t\le1}|R_\alpha(t)-t^\alpha|\le\tau.
\]
The term count is $O_\alpha((1+\log(1/\tau))^2)$. Construction time,
coefficient bits, and the expanded dense numerator and denominator
sizes are polynomial in the tolerance encoding and accuracy bits.
\end{lemma}
\begin{proof}
For $t\ge0$ set
\[
 I_\alpha(t)=\int_0^\infty\frac{t s^{\alpha-1}}{t+s}\,ds.
\]
It converges for $t>0$, and its integrand is zero for $t=0$.
Substitution $s=tu$ for $t>0$ proves
$I_\alpha(t)=t^\alpha I_\alpha(1)$, also valid at zero. For $0\le t\le1$,
truncating to $[2^{-L},2^L]$ loses at most
\[
 \alpha^{-1}2^{-\alpha L}
       +(1-\alpha)^{-1}2^{-(1-\alpha)L}.
\]
On panel $[2^k,2^{k+1}]$, substitute $s=2^ku$ and use the $n$-point
Gauss--Legendre rule on $[1,2]$, with nodes $u_j\in(1,2)$ and positive
weights $w_j$ summing to one. It is exact through degree $2n-1$.
The resulting approximation is
\[
 Q(t)=\sum_{k=-L}^{L-1}\sum_{j=1}^n
       w_j2^{k\alpha}u_j^{\alpha-1}
       \frac{t}{t+2^ku_j}.
\]
For completeness, the error is uniform even at zero. The panel integrand
\[
 F_{k,t}(u)=2^{k\alpha}u^{\alpha-1}\frac{t}{t+2^ku}
\]
is analytic on a neighborhood of $|u-3/2|\le1$. On this disk
$\mathop{\rm Re}u\ge1/2$, so the principal power is analytic,
$|u^{\alpha-1}|\le2$, and $|t/(t+2^ku)|\le1$. Cauchy's coefficient
bound makes the degree-$(2n-1)$ Taylor remainder on $[1,2]$ at most
$4\,2^{k\alpha}2^{-2n}$. Exactness and positivity of the quadrature
bound its integration error by twice that remainder. Thus
\begin{equation}
 \sup_{[0,1]}|I_\alpha-Q|
 \le \alpha^{-1}2^{-\alpha L}
       +(1-\alpha)^{-1}2^{-(1-\alpha)L}
       +16L\,2^{L-2n}.
 \label{eq:a-law-quadrature}
\end{equation}
For target unnormalized error $\zeta$, take $L=O_\alpha(1+\log(1/\zeta))$
and $n=O_\alpha(1+\log(1/\zeta))$ so the right side is at most
$\zeta/2$. This gives $J=2Ln$ terms. Integer choices can be made from
rational upper bounds on the displayed powers, or from these fixed
constant inequalities with additional slack.

Here is an effective rounding argument. The Legendre recurrence
constructs a rational polynomial of degree $n$ and polynomial coefficient
bits. Root isolation and the positive weight formula
$2/((1-v_j^2)P'_n(v_j)^2)$ on $[-1,1]$, followed by the affine change
of interval, give rational enclosures for each node and weight.
Their fixed rational powers are obtained by isolating the corresponding
positive algebraic roots. Each individual number has polynomial algebraic
degree and encoding length; one never forms their joint algebraic field.
Polynomial-time arbitrary-precision enclosures also follow directly from
\citet[Sections 1 and 7.2]{johansson2018-gauss-legendre}.

Before rounding, $s=2^ku_j\ge2^{-L}$ and
$0<a=w_j2^{k\alpha}u_j^{\alpha-1}\le2^L$. Enclose and round each $a,s$
upward to positive dyadic rationals with error at most $h\le1$.
This works even for an arbitrarily small weight: an upper endpoint of
a sufficiently narrow enclosure is positive. Along these rounding
segments, $s\ge2^{-L}$ and $a\le2^{L+1}$. For every $t\ge0$,
\[
 \left|\frac{\partial}{\partial a}\frac{at}{t+s}\right|\le1,
 \qquad
 \left|\frac{\partial}{\partial s}\frac{at}{t+s}\right|
       =\frac{at}{(t+s)^2}\le\frac as\le2^{2L+1}.
\]
Choose a dyadic $h$ with
$Jh(1+2^{2L+1})\le\zeta/2$. Its required precision is
$O(\log(1/\zeta)+L+\log J)$ bits, and the rounded sum $Q_{\Q}$ has
uniform error at most $\zeta$. Its coefficients remain positive.

On $[1,2]$, $s^{\alpha-1}/(1+s)\ge1/6$, so
$I_\alpha(1)\ge1/6$. If $\zeta\le1/12$, then $Q_{\Q}(1)\ge1/12$.
Set $R_\alpha(t)=Q_{\Q}(t)/Q_{\Q}(1)$. Since $t^\alpha\le1$,
\[
 |R_\alpha(t)-t^\alpha|
 \le\frac{|Q_{\Q}(t)-I_\alpha(t)|+
             t^\alpha|I_\alpha(1)-Q_{\Q}(1)|}{Q_{\Q}(1)}
 \le24\zeta.
\]
Take $\zeta\le\min\{1/12,\tau/24\}$. Normalization is exact rational
arithmetic. Products of the $J$ positive linear denominators have degree
$J$ and polynomial coefficient bits: denominator bit lengths add, and
coefficient growth adds at most polynomially many further bits.
Summing the corresponding numerators has the same bounds. This proves
the dense encoding and construction claims.
\end{proof}

\begin{lemma}[Rational constitutive approximation]
\label{lem:a-law-constitutive}
Fix a nonempty finite family $\mathcal Q\subset\Q\cap(1,\infty)$. Given rational
$B>0$ and $\delta>0$, one can construct for every $q\in\mathcal Q$ a
rational function $\psi_q$ which is smooth, strictly increasing on $\R$,
vanishes at zero, and has a denominator positive on all of $\R$, such that
\[
 \sup_{|x|\le B}|\psi_q(x)-\phi_q(x)|\le\delta.
\]
Its expanded dense encoding and construction time are polynomial in the
input and accuracy bits, with constants depending on $\mathcal Q$.
\end{lemma}
\begin{proof}
Odd integer $q$ use $\psi_q(x)=x^q$ exactly. Otherwise write uniquely
\[
 q=2j+1+2\alpha,\qquad j\ge0\text{ an integer},\quad0<\alpha<1.
\]
Choose $T=2^h\ge\max\{1,B\}$ with $h$ a nonnegative multiple of the
denominator of $2\alpha$. The rational $T^{2\alpha}$ has polynomial
bit length because the exponent family is fixed. Put
\[
 \psi_q(x)=T^{2\alpha}x^{2j+1}
                 R_\alpha(x^2/T^2).
\]
The corresponding expression with $R_\alpha(t)$ replaced by $t^\alpha$
is exactly $\phi_q(x)$. Apply \cref{lem:a-law-positive-rational} with
$\tau\le\min\{1,\delta/(T^{2\alpha}B^{2j+1})\}$.
Every term of $\psi_q$ is a positive rational multiple of
$x^{2j+3}/(x^2+c)$ for a rational $c>0$, and
\[
 \frac{d}{dx}\frac{x^{2j+3}}{x^2+c}
 =\frac{x^{2j+2}((2j+1)x^2+(2j+3)c)}{(x^2+c)^2}.
\]
It is positive for $x\ne0$ and nonnegative at zero. Therefore the sum
is strictly increasing, is odd, vanishes at zero, and is unbounded in
both directions. Its denominator is a product of positive quadratics.
The number of terms, their bit lengths, and expansion sizes are polynomial
as in the preceding lemma; $j$ and all rational exponent denominators
are fixed. For $1<q<3$ this is the case $j=0$. The extra odd power proves
the same assertion for fixed exponents above three.
\end{proof}

\begin{lemma}[Rational-law local optimization]
\label{lem:a-law-rational-core}
Suppose $g_e(x)=\beta_e\psi_e(x)$, where the densely encoded rational
$\psi_e=P_e/Q_e$ are strictly increasing, vanish at zero, and
$Q_e(x)>0$ for all real $x$. Resistances have independent positive
rational intervals. At fixed block rank, the face construction and
exact local algebraic optimization of \cref{thm:a-law-polynomial}
remain polynomial in the rational encoding, provided polynomial-bit
rational bounds on the laws on $[-B,B]$ are supplied. Thus terminal-drop optima have arbitrarily accurate rational intervals;
each nomination face has separate exact local algebraic optimizers, and
global face selection is additive. Arc-flow optima have exact local
algebraic values and scenarios.
\end{lemma}
\begin{proof}
The face reduction is \cref{lem:a-law-faces}. In a local core $z$,
substitute $X_e(z)$ into $P_e,Q_e$, and form the positive product
$Q_*(z)=\prod_e Q_e(X_e(z))$. Multiplying every cycle equation by $Q_*$
gives polynomial coefficients of the resistance leaves and introduces
no spurious roots. The total degree is at most the sum of the substituted
degrees. In a fixed number of variables, expansion has polynomially
many monomials and polynomial coefficient bits.

For a path-drop objective introduce one extra core coordinate $v$ and
the equation
\[
 \sum_e\omega_{H,e}\beta_e P_e(X_e)
           \prod_{f\ne e}Q_f(X_f)-vQ_*(z)=0.
\]
Bound $v$ by the supplied rational path-drop bound. Treat this as one
extra linking equation and maximize the core-only objective $v$ using
\cref{lem:a-law-dense}. The core dimension and equation count each
increase by just one. For an arc use the affine objective $X_e(z)$
without that extra equation. Keep independent block values separate
and sum rational intervals for the pressure objective as before.
\end{proof}

\subsection{Physical error and rational scenario recovery}
Fix $\mathcal Q\subset\Q\cap(1,\infty)$ nonempty and finite. Throughout this subsection
let
\begin{align*}
 H&=\left\lceil\max\mathcal Q\right\rceil,&
 c_*&=2^{1-H},& T_0&=\max\{1,B\},\\
 M&=H T_0^{H-1},& N_0&=B T_0^{H-1}.
\end{align*}
Thus $H$ is a fixed integer, $c_*>0$ is rational, and on $[-B,B]$ every
$\phi_q$ and its derivative are bounded in absolute value by $N_0$ and
$M$, respectively. These bounds have polynomial rational bit length.

\begin{lemma}[Constitutive error controls the physical state]
\label{lem:a-law-state-error}
Let $m>0$, $\beta_L=\min_e L_e$, and $\beta_U=\max_e U_e$. Suppose
$|\psi_e-\phi_{q_e}|\le\delta$ on $[-B,B]$, where the $\psi_e$ are
continuous strictly increasing laws vanishing at zero. For the same
admissible $(b,\beta)$, let $x,\pi$ be the original state and $y,\widehat\pi$
the state with laws $\beta_e\psi_e$. For $0<\eta\le1$, the condition
\begin{equation}
 \delta\le\frac{c_*\beta_L\eta^H}{m\beta_U}
 \label{eq:a-law-delta-budget}
\end{equation}
implies
\begin{equation}
 \|x-y\|_\infty\le\eta,\qquad
 |(\pi_s-\pi_t)-(\widehat\pi_s-\widehat\pi_t)|
       \le m\beta_U(M\eta+\delta).
 \label{eq:a-law-state-error}
\end{equation}
Both estimates are uniform on the full scenario domain.
\end{lemma}
\begin{proof}
For $q>1$ the signed power satisfies
\begin{equation}
 (\phi_q(a)-\phi_q(b))(a-b)
              \ge2^{1-q}|a-b|^{q+1}
              \ge c_*|a-b|^{q+1}\quad(q\in\mathcal Q).
 \label{eq:a-law-strong}
\end{equation}
For equal signs assume $a=b+d\ge b\ge0$ after symmetry. The function
$(b+d)^q-b^q$ is increasing in $b$, so it is at least $d^q$.
For opposite signs convexity gives
$|a|^q+|b|^q\ge2^{1-q}(|a|+|b|)^q$.
Multiplying these increment bounds by $|a-b|$ proves the inequality.

Both flows lie in $[-B,B]^E$ by acyclicity, and $A(x-y)=0$.
Pairing this circulation with the difference of the two potential
gradients yields
\[
 \sum_e\beta_e[\phi_{q_e}(x_e)-\phi_{q_e}(y_e)](x_e-y_e)
 =\sum_e\beta_e[\psi_e(y_e)-\phi_{q_e}(y_e)](x_e-y_e).
\]
Let $D=\|x-y\|_\infty$. If $D=0$ the flow bound is immediate. Otherwise
select an edge $e$ attaining $D$. Nonnegativity of all terms on the
left, \cref{eq:a-law-strong}, and the uniform error on the right give
$c_*\beta_L D^{q_e+1}\le m\beta_U\delta D$, hence
$D^{q_e}\le\eta^H$ under \cref{eq:a-law-delta-budget}.
If $D>\eta$ and $D\le1$, then $D^{q_e}\ge D^H>\eta^H$;
if $D>1$, then $D^{q_e}>1\ge\eta^H$. Both contradict the bound.
Thus $D\le\eta$. Each edge-drop error is at most
$\beta_U(M\eta+\delta)$. Summing on a simple terminal path proves the
second estimate.
\end{proof}

\begin{lemma}[Original-law input rounding]
\label{lem:a-law-fractional-rounding}
For two admissible scenarios $(b,\beta)$ and $(b',\beta')$, original-law
terminal drops satisfy
\begin{equation}
 |F_{s,t}(b,\beta)-F_{s,t}(b',\beta')|
 \le M\sum_{e\in P_{st}}U_e\,\|b-b'\|_1
                     +N_0\|\beta-\beta'\|_1.
 \label{eq:a-law-original-lipschitz}
\end{equation}
For a target arc $e=(u,v)$ with original flows $x_e,x'_e$, write the
corresponding endpoint drops as $\Delta,\Delta'$. Then
\begin{equation}
 |x_e-x'_e|^{q_e}\le
 \frac{|\Delta-\Delta'|+N_0|\beta_e-\beta'_e|}{c_*L_e}.
 \label{eq:a-law-inverse-rounding}
\end{equation}
In particular, a right side at most $\xi^H$, $0<\xi\le1$, guarantees
flow error at most $\xi$. Polynomially many input precision bits suffice
for any prescribed pressure or flow accuracy.
\end{lemma}
\begin{proof}
The original laws are continuously differentiable. Add $\rho x$ to each
complete law and use the ordinary electrical adjoint. Its nomination
path bound is $\sum_{e\in P_{st}}(U_eM+\rho)$. At fixed nomination,
resistance differentiation gives
$\partial F_\rho/\partial\beta_e=\jmath_e\phi_{q_e}(x_e)$, because the
added linear term is independent of $\beta_e$. This is bounded by
$N_0$. Integrate over the two box segments and pass to the uniform
zero-smoothing limit as before. This proves
\cref{eq:a-law-original-lipschitz}.

Subtracting $\Delta=\beta_e\phi_{q_e}(x_e)$ and
$\Delta'=\beta'_e\phi_{q_e}(x'_e)$ gives
\[
 \beta_e[\phi_{q_e}(x_e)-\phi_{q_e}(x'_e)]
       =\Delta-\Delta'+(\beta'_e-\beta_e)\phi_{q_e}(x'_e).
\]
Use the increment version of \cref{eq:a-law-strong}, the bound $N_0$,
and $\beta_e\ge L_e$. The final implication follows from $q_e\le H$
exactly as in \cref{lem:a-law-state-error}. For the endpoint drop, take
the one-edge path in \cref{eq:a-law-original-lipschitz}; its nomination
constant is $U_eM$. Thus a desired inverse error requires only a fixed
power $\xi^H$ of the desired flow error. All constants have polynomial
rational bit length.
\end{proof}

\begin{theorem}[Fixed fractional powers: polynomial accuracy-bit algorithms]
\label{thm:a-law-fractional}
Fix a block-rank bound $\mathfrak r$ and a nonempty finite family
$\mathcal Q\subset\Q\cap(1,\infty)$. On each edge choose $q_e\in\mathcal Q$
and the law $\beta_e\phi_{q_e}$, with independent positive rational
resistance intervals. For a nonempty balanced rational nomination box,
one can compute a rational interval of width at most $\eps$ containing
the maximum prescribed terminal drop, and rational admissible nomination
and resistance inputs within $\eps$ of that maximum. The same guarantees
hold for the maximum and minimum signed flow on any prescribed arc.
Running time is polynomial in input size and accuracy bits, with its
exponent and constants allowed to depend on $\mathfrak r,\mathcal Q$.
No exact threshold comparison or rational physical-state output is claimed.
\end{theorem}
\begin{proof}
Replace $\eps$ by $\min\{\eps,1\}$. The edgeless and $B=0$ cases are
immediate. Put $a=\eps/8$ and choose rational
\[
 \eta=\min\{1,a,a/(2m\beta_U M)\},\qquad
 \delta=\min\{1,a/(2m\beta_U),
                     c_*\beta_L\eta^H/(m\beta_U)\}.
\]
Construct the rational $\psi_{q_e}$ by
\cref{lem:a-law-constitutive}. Their flow objectives differ from the
original ones by at most $a$, and their terminal-drop objectives by at
most $a$, uniformly over the scenario domain, by
\cref{lem:a-law-state-error}. A rational absolute drop bound is
$m\beta_U(N_0+\delta)$, so \cref{lem:a-law-rational-core} applies.
All construction degrees and bit lengths are polynomial in the original
input and accuracy bits; $H$ and exponent denominators are fixed.

Compute each surrogate block optimum exactly, retaining its local
algebraic scenario. Enclose each block value to width
$\eps/(8n)$, sum by face, and select the face with greatest lower
endpoint. This gives an interval $[L,U]$ of width at most $\eps/8$
containing the surrogate maximum, and a separately encoded algebraic
scenario whose surrogate value loses at most $\eps/8$. For an arc there
is only one relevant block, and exact local comparison is possible;
one may use the same interval budget. A minimum is obtained by maximizing
the negative objective. The interval $[L-a,U+a]$ contains the original
maximum and has width at most $3\eps/8$.

Apply \cref{lem:a-law-fractional-rounding} to recover rational inputs
for the \emph{original} laws.
For pressure, choose sufficiently narrow nomination and resistance
isolating boxes that \cref{eq:a-law-original-lipschitz} is at most
$\eps/8$. For an arc, use \cref{eq:a-law-inverse-rounding} and require
its right side to be at most $(\eps/8)^H$. Nomination-box intersection
with its rational face and balance contains the exact algebraic point;
rational LP returns an exactly feasible rational point. Resistance
coordinates are rounded in their original intervals independently.
All needed enclosures have polynomial precision. On a terminal core,
round its aggregated nomination first and then greedily disaggregate;
its objective is determined by that aggregate, so the sensitivity bound
is applied on that core. Complete irrelevant coefficients with rational
endpoints. No exact original-law flow computation is needed for this
recovery.

The output's original objective loses at most $a$ in comparing the two
optima, $\eps/8$ in face selection, $a$ in comparing the two laws at the
selected algebraic scenario, and $\eps/8$ in rounding. This totals at
most $\eps/2$, within the requested tolerance. In particular, for arc
optimization one starts from an arc optimizer, not a pressure optimizer
with possibly different target resistance.
\end{proof}

The preceding theorem includes any fixed finite family in $(1,3)$,
and the odd-power factor in \cref{lem:a-law-constitutive} extends it to
all fixed rational exponents greater than one. The approximation
principle itself is not a novelty claim: the contribution is its use
with the block-rank reduction, rational encoding, uniform physical-error
transfer, and rational scenario recovery. If exponents are arbitrary
binary input, $\alpha$ or $1-\alpha$ can be exponentially small and the
present dyadic construction can require exponentially many panels.
Also the odd power and the exponent $H$ would no longer be fixed.
No algorithm polynomial in such an exponent encoding is asserted.
Exact integer power laws have the polynomial representations already
treated above, including the linear law at the endpoint $q=1$.
\Cref{thm:a-law-srs} remains compatible with these additive guarantees:
an accuracy interval need not decide an exact equality or a threshold
closer than the requested error.

For example, the water head-loss exponent $1.852$ in
\citet[Example 2.3, May 2018 manuscript]{gro2017-algorithmic-results-for-potential-based}
is the fixed reduced rational $463/250$. Its even denominator puts that
mathematical law within both the single-cycle exact barrier and the
fixed-family additive theorem.
